% tilics/booklet/main.tex

% 2019-05-11 christian.tschudin@unibas.ch

\documentclass{tilics}

% \minibooklettrue

\ifminibooklet
  \newcommand{\artcount}{4}
  \setlength{\pdelta}{4.6mm}
\else
  \newcommand{\artcount}{5}
  \setlength{\pdelta}{2.3mm}
\fi


% ----------------------------------------------------------------------
\begin{document}
  % \pagecolor{yellow}
  \pagecolor{UnibasMint}
  \hfill\includegraphics{figs/019-turing-machine.pdf}

  \vspace*{\fill}{\color{UnibasRed}
    \noindent{\Huge\bf 2\raisebox{0.5em}{\LARGE\artcount} Things I Learned}
    \vskip 1mm
    \noindent{\Huge\bf in Computer Science}
  }

  \vskip 4mm\noindent{\large
    University of Basel, Dept of Mathematics and Computer Science
  }

  % 2nd page (even) with affiliation
  \newpage\pagecolor{white}\rightBar

  \noindent
  {\small\copyright\ 2019
  University of Basel, Department of Mathematics and Computer Science}

  \vspace*{-2pt}\noindent
  {\small Spiegelgasse 1, CH -- 4051 Basel, Switzerland}
  
  \vspace*{-3pt}\noindent
  {\small\tt https://dmi.unibas.ch/}
  
  \vspace*{\fill}
  \noindent
  \includegraphics{figs/CC-by-sa.pdf}\hspace*{0.4em}
  \raisebox{1.5em}{\begin{minipage}{8cm}\small\raggedright
    This work is licensed under a {\em Creative Commons Attribution-ShareAlike 4.0 International License}\\
    (CC BY-SA 4.0). 
  \end{minipage}}

  \vspace*{-0.4em}\noindent
 {\small\tt https://creativecommons.org/licenses/by-sa/4.0/legalcode}
  % {\small University of Basel, Dept of Mathematics and Computer Science}

  % 3rd page (odd) with (?author and) title
  \newpage
  \newgeometry{left=2cm,right=2cm}
  \voffset 0.4in%\hoffset 0.1in

  \leftBar
  \vspace*{\fill}
  \noindent\hfil{\Huge\bf 2\raisebox{0.5em}{\LARGE\artcount} Things I Learned}
  \vskip 1mm
  \noindent\hfil{\Huge\bf in Computer Science}

  \vskip 8pt
  \noindent\hfil{\large Department of Mathematics and Computer Science}
  \vskip 30mm
  \noindent\hfil\includegraphics{figs/UniBas_Logo_EN_Schwarz_Trans_RGB_65.pdf}
  \vspace*{-8mm}
 
  
  % 4th page (even) with Author's Note, a
  \newpage\rightBar %\vspace*{\fill}

  \vspace*{1cm}
  \noindent{\Large Author's Note}
  \vskip 5mm\small
  \noindent
Studying Computer Science is like taking a deep dive into an endless
ocean of ideas, concepts, and tools. It is a vast ecosystem where
everything is interconnected, where exploring one thing opens doors to
a plethora of other, equally interesting, things. It is easy to get
lost in this wide landscape, even more so if it is not one's main
field of expertise.

This is why Computer Science students of the University of Basel set
out to identify the essential concepts and ideas they discovered
during their studies. The intention was to go beyond the scope of
traditional textbooks and also include the pecularities that are not
formally part of a usual Computer Science curriculum. Finally, the
goal was to present the fruit of this work to people outside of our
field.

  % 5th page (odd) with Author's Note, b
  \newpage
  \leftBar
\vspace*{2cm}
The creation of this booklet was organized in the form of a Master's
level seminar, which had its first iteration during the spring
semester 2019. We have listed the participants on the inside back cover.

The result is
{\bf 2\raisebox{0.3em}{\hspace{1pt}\artcount} Things I Learned in Computer Science},
a collection of short illustrated articles. It is available in the
form of this small booklet, and online on the companion website
$[$1$]$. Due to the success and the positive feedback, we will continue
this work and are now planning for a next edition of the seminar.

\vskip 0.5em\noindent{\raggedleft\hfill
We hope that you enjoy these info bytes\\
\hfill and that you pass them on to your friends!}

{\raggedright
\vspace*{\fill}\noindent
{\small\tt $[$1$]$ https://tilics.dmi.unibas.ch/}}

\vspace*{\fill}

%  % 7th page (odd) with continuation of Author's Note
%  \newpage\leftBar\vspace*{\fill}
%  \noindent{\small\blindtext}
%  \vspace*{\fill}

  \newcommand{\tilicsArtNum}{???}
  
  % loop over all article
  \foreach \n in {0,...,31}{
     \IfFileExists{articles/\artcount .\n/fig.pdf}{
       \typeout{Article \n  ...}

       % left side (even page):
       \inputFig{articles/\artcount .\n/fig.pdf}

       % right side (odd page):
       \input{articles/\artcount .\n/num.tex}
       \inputText{articles/\artcount .\n/txt.tex}
     }{}
  }

  % index page
  \newpage\rightBar

  \vspace*{-0.5em}
  \noindent
  {\Large\bf List of topics}

  \vskip 2em
  \noindent
  \ifminibooklet
    \begin{tabular}{rp{1em}lp{3em}rp{1em}l}
    0 && Hello World           &&    8 &&  Monte Carlo method\\
    1 && Automation            &&    9 && System vs service uptime \\
    2 && Binary system         &&    10 && Camel case \\
    3 && Finite state machine  &&    11 && Programming bugs \\
    4 && Turing machine        &&    12 && Version control systems \\
    5 && Polygon mesh          &&    13 && Deadlock \\
    6 && XOR                   &&    14 && Year 2038 problem \\
    7 && Zero vs Null          &&    15 && Reverse engineering \\
    \end{tabular}
  \else
    \begin{tabular}{rp{1em}lp{3em}rp{1em}l}
    0 && Hello World           &&    16 && Programming bugs \\
    1 && Handshake             &&    17 && Patching \\
    2 && Automation            &&    18 && Octal system \\
    3 && Binary system         &&    19 && 25\_dec = 31\_oct \\
    4 && Run time complexity   &&    20 && Version control systems \\
    5 && Finite state machine  &&    21 && Interface in software engineering \\
    6 && Turing machine        &&    22 && Deadlock \\
    7 && Polygon mesh          &&    23 && Numeric precision \\
    8 && XOR                   &&    24 && Mutex \\
    9 && Monte Carlo method    &&    25 && Quine \\
    10 && Turing test          &&    26 && Greedy algorithms \\
    11 && ASCII                &&    27 && Speculative execution \\
    12 && System vs service uptime &&    28 && Reverse engineering \\
    13 && Overflow             &&    29 && Brute force attack \\
    14 && Camel case           &&    30 && Year 2038 problem \\
    15 && Zero vs Null         &&    31 && Forking open source projects\\
    \end{tabular}
    \vspace*{-1.5em}
  \fi

  % contributors
  \newpage\leftBar%\vspace*{\fill}

  \vspace*{-0.5em}
  \noindent{\Large\bf Seminar participants (spring semester 2019)}
  \normalsize

  \vskip 1.5em
  \noindent\small
  \begin{tabular}{p{3.8cm}p{3.5cm}p{3.5cm}}
    Fabricio Arend Torres &  Patrick Kahr         &  Christian Tschudin   \\
    Nils Bühlmann         &  Marcel Lüthi         &  Marco Vogt           \\ 
    Rebecca Dold          &  Sebastian Philipps   &  Moira Zuber          \\
    Simon Dold            &  Linard Schwendener   &  \\
    Omnia Kahla           &  Tim Steindl          &  \\
  \end{tabular}
  \vspace*{\fill}%

  % back page
  \newpage\pagecolor{UnibasMint}

  \vspace*{\fill}
  \vspace*{72mm}
  \noindent{\small\tt https://tilics.dmi.unibas.ch/}

\end{document}

% eof
